\section{Further Allocation and Budget Comparisons}
\label{app:finite-budget}
The cost term prices the next completed reward evaluation. Let
$U_a=Q_a+c\sqrt{\log(N+1)/N_a}$ and $\gamma=\beta/\cref$ at a fixed decision state.
Write $b_a$ for the expected completion bill under the declared continuation law.
Equation~\eqref{eq:completion-threshold} gives the exact exchange condition.
With equal optimistic quality and $\gamma>0$, the rule chooses a cheapest
completion. With unequal quality, that condition specifies the saving required to
offset the quality disadvantage. This is a resource-priced allocation rule;
$U_a$ retains its tree-exploration meaning and need not be a calibrated
confidence bound on a stationary arm mean.


Proposition~\ref{prop:priced-tradeoff} compares the priced choice $a$ with
the ordinary-UCT choice $b$ at the same initialized state. If
$|\widehat b_j-b_j|\le\epsilon_j$ for $j=a,b$, pricing bounds the sacrifice
in optimistic quality and a predicted saving greater than
$\epsilon_a+\epsilon_b$ certifies a strict reduction in expected physical cost.
The statement is local to one common tree and cache state. Realized stochastic
bills need not follow their expectations on every call.

For a net-value interpretation, let $v_j$ be the actual expected reward of the
declared completion, and additionally require $|U_j-v_j|\le s_j$.
Writing $M=(U_a-\gamma\widehat b_a)-(U_b-\gamma\widehat b_b)$ gives
\begin{equation}
 (v_a-\gamma b_a)-(v_b-\gamma b_b)
 \ge M-s_a-s_b-\gamma(\epsilon_a+\epsilon_b).
 \label{eq:net-certificate}
\end{equation}
A positive right side certifies a net-value improvement. The current tree
bonus and measured prediction MAE do not supply these error bounds. We use
Equation~\eqref{eq:priced-tradeoff} to audit the implemented rule and reserve
the physical and net certificates for settings with valid error bounds.

\begin{proposition}[Forecast error and allocation]
\label{prop:allocation-error}
Let $a$ maximize $U-\gamma\widehat b$, let $a^*$ maximize $U-\gamma b$ at the
same state, and suppose $|\widehat b_j-b_j|\le\epsilon_j$ for these two actions.
For $\gamma\ge0$,
\begin{equation}
 (U_{a^*}-\gamma b_{a^*})-(U_a-\gamma b_a)
 \le\gamma(\epsilon_{a^*}+\epsilon_a).
 \label{eq:allocation-error}
\end{equation}
\end{proposition}
The inequality follows by inserting the two estimated indices and using
their ordering. Summing it bounds cumulative loss against the true-index
choice at each visited state by the corresponding sum of error radii.
The comparator in this statement uses the same observed state. Comparisons
between whole search policies must also account for their different histories
and caches.

Two consequences explain the choice of horizon and cache information. First,
suppose both candidates have the same next-edge bill $e$, with remaining tails
$d_a,d_b$. If $0<U_b-U_a<\gamma(d_b-d_a)$, next-edge chooses $b$, while the full
bill chooses $a$. The absolute truncation error of $e$ is exactly $d_j$.
Second, scale every cold bill in the two-cache witness of
Proposition~\ref{prop:graph-separation} by $H>0$. Any allocation that sees only
the identical prefix totals incurs expected excess cold cost at least $H/2$
in one of the two states. Incidence-aware allocation has zero excess in both.
Thus aggregation and horizon truncation can each erase the cost difference
needed by~\eqref{eq:completion-threshold}.

\begin{proposition}[Coverage under separate cache histories]
\label{prop:budget-coverage}
Consider two search policies with their own realized, nonnegative evaluation
bills. Suppose each of the first $m$ bills of the first policy is at most
$\ell$, and each of the first $k$ bills of the second is at least $h$.
If $m\ell\le B<kh$, every prefix of the first $m$ evaluations fits budget $B$,
whereas the second policy cannot complete its first $k$ evaluations within $B$.
\end{proposition}
Summing the respective bounds proves the claim; nonnegativity bounds each
intermediate prefix cost. Bills include whatever recurring work the declared
budget charges. This conditional statement gives a directly testable benefit:
lower realized completion costs can fund additional reward observations.
It does not order the rewards of those observations. The experiments therefore
report completed evaluations, cost and final quality together, and keep
estimator computation time as a separate measured resource.

\paragraph{Different trajectories.}
Let $A$ and $B$ be the unions of resources materialized by two policies from
the same initial cache $K$. With fixed prices and successful persistent reuse,
their cold-cost difference is exactly
\begin{equation}
 c_w(K,A)-c_w(K,B)
 =c_w(K\cup B,A)-c_w(K\cup A,B).
 \label{eq:exclusive-cost}
\end{equation}
Common materializations cancel; the policies pay differently for resources
exclusive to their own trajectories. Recurring costs add their own difference.
Repeated local savings therefore do not order whole-run expenditure when the
policies explore different workflows. They also do not order endpoint quality.
For example, take current priorities $(U_a,U_b)=(0.6,0.7)$, exact completion
costs $(1,2)$, and $\gamma=0.5$. UCT chooses $b$ and pricing chooses $a$;
both fit a budget of two. If their unknown next rewards are $(0,1)$, pricing
loses both next-reward quality and net value; for $(1,0)$ it wins both. These
two one-decision cases do not assume calibrated $U$ errors. Lean checks the
stated inequalities; they are not a constructed full-tree trajectory theorem.

The finite-budget rationale is consequently precise: use the graph to expose
conditional reuse, estimate the complete reward-producing bill, and trade a
bounded amount of optimistic quality for predicted savings. Evaluate the
resulting policies at a common retained-search-cost ceiling, with the full optimization bill reported separately. The logarithmic theorem below
establishes preservation of long-run quality exploration under persistent cold
costs; it supplies no universal finite-budget advantage over UCT.
